\chapter{Introduction}

\section{Overview}

\VirtualReactor{} is a modular Python framework for the mechanistic simulation
of chemical reactors. Reaction kinetics, reactor models, thermophysical
properties, and numerical integration are implemented as independent
components that can be combined within a common simulation framework.

The current implementation supports batch and plug-flow reactors, reversible
reaction networks, non-isothermal operation, and HDF5-based storage of
simulation results.


\section{Motivation}

Chemical reactor models couple reaction kinetics with reactor-specific mass
and energy balances. Implementing both within a single model limits the
reusability of reaction mechanisms and complicates systematic variation of
reactor configurations and operating conditions.

\VirtualReactor{} separates the chemical reaction model from the reactor
description. This enables the same reaction mechanism to be evaluated in
different reactor configurations and provides a common basis for systematic
parameter studies and optimization.


\section{Design Philosophy}

The framework follows a modular object-oriented design. A
\texttt{Mechanism} describes reaction kinetics and thermodynamics, while
reactor-specific behavior and thermophysical properties are represented by
independent models. The \texttt{Simulation} class combines these components
and performs the numerical integration.

Simulation output is represented independently by
\texttt{SimulationResult}, allowing numerical results to be analyzed and
stored without dependence on the simulation procedure.


\section{Scope}

The current version focuses on deterministic simulations of homogeneous
chemical reaction networks in batch and steady-state plug-flow reactors.
Temperature-dependent kinetics, reversible reactions, reaction heat, and
external heat transfer are supported.

Detailed fluid dynamics, multiphase systems, and spatially resolved transport
are outside the current scope.

Future development will extend the framework toward systematic data
generation and data-driven reactor optimization, including Bayesian
optimization of operating conditions and reactor parameters.

